% Chapter 28: From Gesture to Glyph
\chapter{From Gesture to Glyph}

\chapterepigraph{%
  Writing did not begin when someone decided to record language.\\
  It began when someone noticed that their hand\\
  left marks that others could follow.
}{Flyxion, \textit{Semantic Infrastructure}}

\sidenote{Connects to\\App.\,F gesture\\manifolds;\\ Ch.\,10--13\\on gesture theory}

Part~IX develops the theory of symbols, alphabets, and compression.
This chapter examines the deepest transition in the history of human
communication: the emergence of glyphs from gestures.

The claim of the Gesture Before Symbol framework (Chapters~10--13) is
that symbols are projections of gesture equivalence classes. Applied to
writing, this becomes: glyphs are fossil gestures — motor trajectories
whose temporal structure has been compressed into a static spatial mark.

\section{How Movement Becomes Mark}

The transition from gesture to mark is a projection in the strict
sense of Appendix~B. The gesture $\gamma : [0,T] \to \mathbb{R}^3$
is a trajectory through the hand's configuration space. The mark
$m \in \mathbb{R}^2$ is its trace on a surface — a projection from
the three-dimensional, temporal gesture to the two-dimensional, atemporal
imprint.

\begin{definition}{Trace Projection}{trace-proj}
  The \emph{trace projection} of a writing gesture is the map
  $\Pi_\text{trace} : \gesturesp \to \mathcal{M}$, where
  $\mathcal{M}$ is the space of planar curves, defined by
  $\Pi_\text{trace}(\gamma) = \{(\gamma_x(t), \gamma_y(t)) : t \in [0,T]\}$.
  The trace discards timing, pressure variation, and the full
  three-dimensional trajectory, retaining only the planar ink path.
\end{definition}

The fiber $\Pi_\text{trace}^{-1}(m)$ over a single mark $m$ is enormous:
it contains all gestures that produce the same planar trace regardless of
speed, pressure, pen tilt, or three-dimensional hand movement.
What writing preserves is the \emph{shape} of the gesture; what it
discards is the \emph{dynamics}.

\section{The Stabilization of Traces}

A single person making a single mark produces one trace. A culture
making a particular mark repeatedly, across many individuals and
occasions, produces a distribution over traces. The \emph{stable trace}
— the canonical form that the culture recognizes — is an attractor
in the space of traces: the mode of this distribution.

\begin{definition}{Stable Glyph}{stable-glyph}
  A \emph{stable glyph} is an attractor $g^* \in \mathcal{M}$ in the
  trace space under the dynamics of cultural transmission. It satisfies:
  (1) traces produced by competent writers converge to $g^*$, and
  (2) the basin $\mathcal{B}(g^*)$ is large enough to include the
  natural variation of all competent writers.
\end{definition}

A glyph is stable when its basin of attraction in trace space is
large relative to the variation produced by the community. Letters
whose basins are small — easily confused with neighboring attractors —
are unstable and tend to be replaced or merged with neighboring letters
in the course of script evolution.

\section{Why Writing Systems Emerge}

Writing systems do not emerge when humans develop the capacity to make
marks — that capacity is ancient, predating \textit{Homo sapiens}.
They emerge when two additional conditions are met:

\begin{enumerate}
  \item \textbf{Stable attractors:} the culture has developed a set of
        stable glyphs whose basins are reliably accessible to all
        community members.
  \item \textbf{Compositional semantics:} sequences of glyphs are
        assigned meanings that are not arbitrary — the meaning of a
        sequence is computable from the meanings of its parts by
        a systematic rule (syllabic, alphabetic, logographic).
\end{enumerate}

The first condition is a constraint on the accessibility landscape of
trace space: enough stable attractors must exist, spread far enough
apart, to distinguish the basic units of the language. The second
condition is a constraint on the semantic infrastructure: the system
of glyphs must implement a module with a composable interface.

\section{Gesture Fossilization}

Over centuries, glyphs drift from their gestural origins. The Phoenician
letter \textit{aleph} began as a pictograph of an ox head; the Greek
$\Alpha$ and Latin $A$ retain the shape but have lost all connection
to the ox. The gesture that originally drew the ox head is now a
standardized abstract form produced by culturally trained motor programs
entirely unlike the original sketching motion.

\begin{definition}{Gesture Fossilization}{gesture-fossil}
  A glyph undergoes \emph{gesture fossilization} when:
  \begin{enumerate}
    \item The motor program for producing it is standardized to a fixed
          sequence of strokes regardless of any representational intent,
    \item The fiber $\Pi_\text{trace}^{-1}(g)$ has shrunk: fewer
          gesture trajectories are recognized as producing $g$,
    \item The connection to the original representational gesture is
          no longer accessible to competent writers.
  \end{enumerate}
\end{definition}

Fossilization is irreversible under normal conditions: a fossilized glyph
cannot be ``unfossilized'' by historical knowledge alone. The motor
program for writing $A$ has been overwritten by thousands of repetitions;
the original ox-head gesture cannot be recovered by knowing that $A$
derives from aleph.

\begin{exercise}
  The Chinese character for ``sun'' (日) began as a pictograph of a circle
  with a dot (the sun's disk). Trace its evolution through oracle bone,
  bronze, seal, and modern forms. At what point does fossilization occur?
  What is the trace projection of the modern form versus the original?
\end{exercise}

\begin{exercise}
  Arabic script is cursive — letters connect and change shape depending
  on their position in a word (initial, medial, final, isolated).
  Model this as a context-dependent trace projection: the same glyph
  has different attractors depending on its position in the word.
  What does this imply about the fiber structure of Arabic character recognition?
\end{exercise}

\begin{exercise}[title={Open problem}]
  Can writing systems be designed that maximize the stability of their
  glyphs while minimizing motor acquisition cost? Define both quantities
  formally and propose an optimization criterion. What does the
  optimal alphabet look like?
\end{exercise}
